\chapter{Shared Library Modules}
\label{ch:lib}

The units share five modules in \file{lib/src}: \sig{unpack}
(Section~\ref{sec:unpack}), \sig{cls} (Section~\ref{sec:cls}), \sig{norm}
(Section~\ref{sec:norm}), \sig{round} (Section~\ref{sec:round}) and \sig{pack}
(Section~\ref{sec:pack}). Table~\ref{tab:lib_use} shows which unit uses which
module. A user does not instantiate them directly when using a unit, but must
compile them with it (Section~\ref{sec:files}).

\begin{threetable}{Shared Modules and the Units That Use Them}{tab:lib_use}{Module}{Clocked}{Used by}{2.4cm}{1.8cm}
\sig{unpack} & No & \sig{fadd}, \sig{fmul}, \sig{fdiv}, \sig{fmac}, \sig{ffma}, \sig{fpu} \\ \hline
\sig{cls} & No & \sig{norm}, and the \sig{fdiv\_core}, \sig{fmac\_core} and \sig{ffma\_core} datapaths \\ \hline
\sig{norm} & Yes & \sig{fadd}, \sig{fmul}, \sig{fdiv}, \sig{fpu}; \sig{fmac} and \sig{ffma} normalize inside their cores \\ \hline
\sig{round} & No & \sig{pack} \\ \hline
\sig{pack} & Yes & \sig{fadd}, \sig{fmul}, \sig{fdiv}, \sig{fmac}, \sig{ffma}, \sig{fpu} \\
\end{threetable}

\section{unpack}
\label{sec:unpack}

\sig{unpack} splits one operand into sign, exponent and mantissa and classifies
it. It is combinational. For a subnormal operand it outputs the exponent 1 and a
mantissa with a zero hidden bit. Its parameters are \sig{DATA\_W}, \sig{EXP\_W}
and \sig{MAN\_W} (Table~\ref{tab:common_params}), and its signals are listed in
Table~\ref{tab:unpack}.

\begin{sigtable}{\sig{unpack} Signals}{tab:unpack}
data\_i & Input & DATA\_W & Operand in the format of Section~\ref{sec:format}. \\ \hline
nan & Output & 1 & Exponent field all ones and fraction not zero. \\ \hline
infinite & Output & 1 & Exponent field all ones and fraction zero. \\ \hline
zero & Output & 1 & Exponent field and fraction zero. \\ \hline
sign & Output & 1 & Sign bit. \\ \hline
exponent & Output & EXP\_W & Exponent field, or 1 if the exponent field is zero. \\ \hline
mantissa & Output & MAN\_W & Hidden bit (1 if the exponent field is not zero) followed by the fraction. \\
\end{sigtable}

\section{cls}
\label{sec:cls}

\sig{cls} (count leading symbols) counts how many bits, from the most
significant bit down, equal \sig{symbol}. It is combinational and has one
parameter, \sig{DATA\_W}, the input width (default 32). Its signals are listed
in Table~\ref{tab:cls}.

\begin{sigtable}{\sig{cls} Signals}{tab:cls}
data\_in & Input & DATA\_W & Word to scan. \\ \hline
symbol & Input & 1 & Bit value to count. \\ \hline
data\_out & Output & $\lceil\log_2(\mathtt{DATA\_W}+1)\rceil$ & Number of leading bits equal to \sig{symbol}; \sig{DATA\_W} if all bits are. \\
\end{sigtable}

\section{norm}
\label{sec:norm}

\sig{norm} normalizes a mantissa of \sig{IN\_W} bits, whose top bit
has weight 2, and reduces it to $\mathtt{MAN\_W}+3$ bits. Its output
is registered, so it adds one clock cycle.

\begin{itemize}
\item It shifts the mantissa left by its leading-symbol count from \sig{cls},
  limited so that a subnormal result keeps the minimum exponent.
\item If the exponent is below the minimum, it shifts the mantissa right
  instead, by at most $\mathtt{MAN\_W}+2$ positions, and collects the lost bits
  in a sticky bit.
\item The least significant output bit is the sticky bit of all the bits
  dropped by the reduction.
\end{itemize}

Its parameters are \sig{DATA\_W}, \sig{EXP\_W} and \sig{MAN\_W}
(Table~\ref{tab:common_params}) and those of Table~\ref{tab:norm_params}. Its
signals are listed in Table~\ref{tab:norm}.

\begin{threetable}{\sig{norm} Parameters}{tab:norm_params}{Parameter}{Default}{Description}{2.0cm}{2.2cm}
\sig{IN\_W} & $2\cdot\mathtt{MAN\_W}$ & Mantissa input width. \sig{fadd} and \sig{fdiv} use $\mathtt{MAN\_W}+4$. The quotient of \sig{fdiv} is below 2, so \sig{fdiv\_core} gives its weight-1 bit as the top bit and lowers the exponent by 1. \\ \hline
\sig{RSHIFT} & 1 & 1: signed exponent of EXP\_W+2 bits, and results below exponent 1 shifted right to subnormals. 0: unsigned exponent of EXP\_W bits, never below 1, and no right shift, as in \sig{fadd}. \\ \hline
\sig{LZC\_W} & \sig{IN\_W} & Top bits examined by \sig{cls}, which bounds the left shift: 1 in \sig{fmul}, 2 in \sig{fdiv}. \\
\end{threetable}

\begin{sigtable}{\sig{norm} Signals}{tab:norm}
clk & Input & 1 & Clock. \\ \hline
rst & Input & 1 & Asynchronous reset, active high. \\ \hline
valid\_i & Input & 1 & Input data valid. \\ \hline
valid\_o & Output & 1 & \sig{valid\_i} delayed by one clock cycle. \\ \hline
sign\_i & Input & 1 & Sign. \\ \hline
exp\_i & Input & EXP\_W+2, or EXP\_W with \sig{RSHIFT}~=~0 & Exponent. \\ \hline
man\_i & Input & IN\_W & Mantissa. \\ \hline
res\_special\_i & Input & 3 & Special result code, passed to \sig{res\_special\_o}. \\ \hline
res\_special\_o & Output & 3 & Registered \sig{res\_special\_i}. \\ \hline
sign\_o & Output & 1 & Registered sign. \\ \hline
exp\_o & Output & EXP\_W+1, or EXP\_W with \sig{RSHIFT}~=~0 & Normalized exponent, never negative. \\ \hline
man\_o & Output & MAN\_W+3 & Normalized mantissa with sticky bit. \\
\end{sigtable}

\section{round}
\label{sec:round}

\sig{round} rounds in the mode \sig{rnd\_mode} (Table~\ref{tab:rnd_mode}), from
the three bits below the mantissa: guard, round and sticky. It is combinational.
If the rounded mantissa overflows, it shifts it right and increments the exponent.
If the exponent exceeds the largest finite exponent, the output is the infinity
encoding, or the largest finite number when the mode rounds toward zero for
\sig{sign}.

Its parameter \sig{DATA\_W} is the \emph{mantissa} width (default 24), and
\sig{EXP\_W} is the exponent width (default 8). With \sig{INV}~=~1 (default 0) it
also accepts a mantissa in one's complement when \sig{sign} is 1, as \sig{fmac} and
\sig{ffma} produce it. The module always takes three extra bits. Its signals are listed in Table~\ref{tab:round}.

\begin{sigtable}{\sig{round} Signals}{tab:round}
exponent\_i & Input & EXP\_I\_W & Exponent, never negative. \sig{EXP\_I\_W} is EXP\_W+1, or EXP\_W in \sig{fadd}, whose exponent never exceeds $2^{\mathtt{EXP\_W}}-1$ before rounding. \\ \hline
mantissa\_i & Input & DATA\_W+3 & Mantissa followed by the guard, round and sticky bits. \\ \hline
sign & Input & 1 & Sign of the result, for the directed modes. \\ \hline
rnd\_mode & Input & 3 & Rounding mode. \\ \hline
exponent\_o & Output & EXP\_W & Rounded exponent. \\ \hline
mantissa\_o & Output & DATA\_W & Rounded mantissa. \\ \hline
round & Output & 1 & The mantissa was incremented. \\
\end{sigtable}

\section{pack}
\label{sec:pack}

\sig{pack} instantiates \sig{round}, assembles the result word and registers it,
so it adds one clock cycle. If the rounded mantissa has a zero top bit, it writes
a zero exponent field (a subnormal or zero result). A zero from the datapath is an
exact cancellation, so it is $+0$, or $-0$ when \sig{rnd\_mode} is 2; the cores send
zero operands as special zero codes with their sign. A special result code other
than the computed result replaces the result. Its parameters are \sig{DATA\_W},
\sig{EXP\_W} and \sig{MAN\_W} (Table~\ref{tab:common_params}), and \sig{INV},
passed to \sig{round}; \sig{fmac} and \sig{ffma} set it to 1. Its signals are listed in Table~\ref{tab:pack}.

\begin{sigtable}{\sig{pack} Signals}{tab:pack}
clk & Input & 1 & Clock. \\ \hline
rst & Input & 1 & Asynchronous reset, active high. \\ \hline
valid\_i & Input & 1 & Input data valid. \\ \hline
valid\_o & Output & 1 & \sig{valid\_i} delayed by one clock cycle. \\ \hline
rnd\_mode & Input & 3 & Rounding mode, passed to \sig{round}. \\ \hline
res\_special\_i & Input & 3 & Special result code \{kind, sign\}. Kind 0: the computed result; 1: the canonical NaN; 2: an infinity; 3: a zero, with the sign. \\ \hline
sign\_i & Input & 1 & Sign. \\ \hline
exp\_i & Input & EXP\_I\_W & Exponent, never negative; see \sig{round}. \\ \hline
man\_i & Input & MAN\_W+3 & Mantissa with guard, round and sticky bits. \\ \hline
data\_o & Output & DATA\_W & Result word. \\
\end{sigtable}
